\section{Extensions, threshold cases, and the \texorpdfstring{$\SUtwo$}{SU(2)} model}
\label{sec:extensions}

\subsection{Several frustrated staples}

The private-link construction realizes any central staple tuple.  In
particular, take $k$ copies of $z^{-1}$ and $m-k$ copies of $\one$.  The
conditional energy is
\begin{equation}\label{eq:k-energy}
 F_{z,k}(U)=(m-k)\dist(\one,U)^2+k\dist(z,U)^2.
\end{equation}

\begin{proposition}[The $k$-versus-$(m-k)$ competition]
\label{prop:k-frustrated}
For $0\leq k\leq m$,
\begin{equation}\label{eq:k-minimum}
 \min_UF_{z,k}(U)=\frac{k(m-k)}{m}L^2.
\end{equation}
For $0<k<m$, the minimizers are exactly the points on minimizing geodesics
from $\one$ to $z$ satisfying
\begin{equation}\label{eq:k-distances}
 \dist(\one,U)=\frac{k}{m}L,
 \qquad
 \dist(z,U)=\frac{m-k}{m}L.
\end{equation}
For $k=0$ and $k=m$, the unique minimizers are $\one$ and $z$,
respectively.
\end{proposition}

\begin{proof}
For $a=\dist(\one,U)$ and $b=\dist(z,U)$, use $a+b\geq L$ and
\begin{equation}\label{eq:k-square}
 (m-k)a^2+kb^2
 =\frac{k(m-k)}m(a+b)^2
  +\frac1m\bigl((m-k)a-kb\bigr)^2.
\end{equation}
Equality gives \eqref{eq:k-minimum}--\eqref{eq:k-distances}, and points on a
minimizing geodesic realize it.  The endpoint cases are immediate.
\end{proof}

The transformation $k\leftrightarrow m-k$ exchanges the roles of $\one$ and
$z$.  At $k=m/2$ the minimizers are geodesic midpoints.  Nonuniqueness can
come from multiple minimizing geodesics, not from the scalar optimization.

\subsection{Threshold behavior}

For one frustrated staple the target distance on the minimizing set is
$b_*=(m-1)L/m$.  The three radius regimes are:
\begin{center}
\begin{tabular}{@{}lll@{}}
\toprule
Radius & Exponential conclusion & Probability conclusion \\
\midrule
$0<\rho<b_*$ & rate $-m(b_*-\rho)^2$ & good probability $\to0$ \\
$\rho=b_*$ & rate $0$ & no limit asserted here \\
$\rho>b_*$ & all minimizers are good & good probability $\to1$ \\
\bottomrule
\end{tabular}
\end{center}
At the equality threshold, the strict good set has the minimizing set on its
boundary.  A logarithmic large-deviation principle alone does not determine
the limiting fraction of mass, so we make no stronger statement.

\subsection{The \texorpdfstring{$\SUtwo$}{SU(2)} class-angle formula}

Normalize the bi-invariant metric on $\SUtwo$ so that every element can be
written
\[
 U=\cos\theta\,\one+i\sin\theta\,\boldsymbol n\cdot\boldsymbol\sigma,
 \qquad 0\leq\theta\leq\pi,
\]
with $\dist(\one,U)=\theta$.  For $z=-\one$, $L=\pi$ and
$\dist(z,U)=\pi-\theta$.  Radial Haar integration is
\begin{equation}\label{eq:radial-haar}
 \int_{\SUtwo}f(\theta(U))\,\Haar
 =\frac2\pi\int_0^\pi f(\theta)\sin^2\theta\,\dd\theta.
\end{equation}
Therefore the one-frustrated conditional class-angle density is
\begin{equation}\label{eq:su2-radial-density}
 \gamma_{t,m}^{-}(\dd\theta)
 =\widetilde Z_{t,m}^{-1}\sin^2\theta\,
 p_t(\theta)^{m-1}p_t(\pi-\theta)\,\dd\theta,
\end{equation}
with energy
\begin{equation}\label{eq:su2-rate}
 I_m^{\mathrm{raw}}(\theta)
 =(m-1)\theta^2+(\pi-\theta)^2
 =m\left(\theta-\frac\pi m\right)^2
  +\frac{m-1}{m}\pi^2.
\end{equation}
The Haar zero at $0$ and $\pi$ is subexponential and does not change the
rate.  Notice that the class angle has a unique minimizer $\pi/m$, while the
corresponding group minimizers form a conjugacy class when $0<\pi/m<\pi$.

In four dimensions $m=6$.  The selected link concentrates at $\theta=\pi/6$
and the frustrated target at $\pi-\theta=5\pi/6$.  For
$0<\rho<5\pi/6$,
\begin{equation}\label{eq:su2-four-dimensional-rate}
 \lim_{t\downarrow0}2t\log\gamma_{t,6}^{-}
 \{\pi-\theta<\rho\}
 =-6\left(\frac{5\pi}{6}-\rho\right)^2.
\end{equation}

\subsection{Low dimensions and trivial center}

The geometry has no exceptional low-dimensional case once $d\geq2$:
\begin{center}
\begin{tabular}{@{}cccc@{}}
\toprule
$d$ & $m=2(d-1)$ & selected-link location & private opposites \\
\midrule
$2$ & $2$ & $(0,1)\to(1,1)$ & $2$ \\
$3$ & $4$ & $(0,1,1)\to(1,1,1)$ & $4$ \\
$4$ & $6$ & $(0,1,1,1)\to(1,1,1,1)$ & $6$ \\
\bottomrule
\end{tabular}
\end{center}
Dimension one is excluded because there are no plaquettes.

If $\cZ(G)$ is trivial, Proposition~\ref{prop:central-surjectivity} supplies
no nonidentity $z$ and the theorem is silent.  Noncentral frustration may
still affect conditional laws, but independent staple prescription then
interacts with ordering and conjugation; no such extension is claimed.  The
paper therefore does not assert a common obstruction for every compact
simple gauge group.
